\subsection{Marked two-primary fixed-root compatibility}

\begin{result}{E.170}{Proved}{An eighth-root mark suffices for the two-primary gluing at primes five modulo eight}
Let \(B\in\mathbb Q^{\times}\) have odd valuation at an odd
prime \(p\), and let \(H\ge3\). Let \(F/\mathbb Q\) be a finite
cyclotomic field containing \(\mu_{2^H}\) and \(\sqrt B\), and assume
\[
 R=F(B^{1/2^H}),\qquad [R:F]=2^{H-1}.
 \tag{E.170.1}
\]
Choose compatible roots. Let \(L/\mathbb Q\) be a finite
Galois field containing \(F\) and \(B^{1/8}\). Suppose that its
additional data consist of finite actual tangential action quotients
of the tripod and finite geometric prime-power section quotients,
with their common arithmetic projection retained. The two-primary
section retains the marked Kummer coordinate of \(B\), which may
also be a specified combination of its abelian coordinates and
cyclotomic Kummer coordinates, at least modulo eight. One may
also include additional Galois fields of odd degree over \(F\).

Let \(g\in\operatorname{Gal}(L/\mathbb Q)\) be an already
compatible prescribed arithmetic and section Frobenius value such
that
\[
 \chi_2(g)\equiv5\pmod8,\qquad g(B^{1/8})=B^{1/8}.
 \tag{E.170.2}
\]
Then \(g\) can be conjugated by an element which is geometrically
two-primary and fixes all the retained arithmetic data and all the
odd-primary geometric data, so that the conjugated value extends to
\(LR\) while fixing a \(2^H\)-th root of \(B\).
Consequently Chebotarev supplies arbitrarily large primes with the
prescribed cyclotomic and section conjugacy data for which \(B\)
is a \(2^H\)-th power in the residue field. Any finite forbidden
set can be avoided.

For an \(A\)-support auxiliary prime, the prescribed cusp-zero
value fixes its marked complementary eighth root. Thus this result
applies to \(B=b/c\) in the coefficient forms of E.154, using
\(p=41\); it applies equally to its inverse. At the infinity
boundary one may instead use \(w=-b/a\), with the mark
\(\kappa(w)=\kappa(b/c)-\kappa(a/c)-\kappa(-1)\), using
inverse compatible roots of \(-1\). Retain \(\mu_{16}\) for
this eighth-root mark. The unit-one infinity cusp fixes it.
It removes the
two-primary fixed-root compatibility issue for such already
compatible Frobenius data. It neither creates those prior data nor
supplies an infinite compatible sequence of sections.
\end{result}
\begin{proof}
We first produce the permitted geometric conjugations without
assuming that the entire geometric image is free or that it equals
its Frattini subgroup. Let \(\rho_\ell\) be the actual tangential
action on the free pro-\(\ell\) group, and set
\[
 K=\ker\left(G_{\mathbb Q}\longrightarrow
                    \prod_\ell\operatorname{Aut}(F_\ell)\right).
 \tag{E.170.3}
\]
Choose complex conjugation \(c\) and tame inertia \(\tau_p\)
whose two-primary tame coordinate is a generator. As proved in
E.127, the commutator \(h=[c,\tau_p]\) has actual tangential
action supported only at \(p\), and that action belongs to a
pro-\(p\) group: it is trivial on the abelianization, and the
corresponding automorphism kernel is pro-\(p\). Here the based
unramifiedness assertion, not merely outer unramifiedness, is
essential \citep[Definition~2.0 and Theorem~2.1, pp.~198--200]
{AuditWojtkowiak2007}. Let \(e_2\in\widehat{\mathbb Z}\) have
two-primary component one and every other component zero. Then
\[
                         k=h^{e_2}\in K.
 \tag{E.170.4}
\]
Its procyclic closure is pro-two. The cyclotomic character of a
commutator is one, so the Kummer cocycle is a homomorphism on that
closure. Its value is
\[
 \kappa_2(B)(k)=-2u\,v_p(B),\qquad u\in\mathbb Z_2^\times.
 \tag{E.170.5}
\]
Indeed the cocycle formula for a commutator is
\(\kappa([c,\tau_p])=(1-\chi(\tau_p))\kappa(c)
 +(\chi(c)-1)\kappa(\tau_p)\); here \(\chi_2(\tau_p)=1\),
\(\chi_2(c)=-1\), and tame Kummer inertia is \(u v_p(B)\).
Taking the two-primary profinite power leaves this value unchanged.

The element \(k\) fixes every root of unity and hence \(F\).
On \(R/F\), formula (E.170.5) is a generator: in the coordinate
\(B^{1/2^H}\mapsto\zeta_{2^H}^{a}B^{1/2^H}\), the group
\(\operatorname{Gal}(R/F)\) consists exactly of the even
residues \(a\), and the displayed value is twice a unit.
On the other hand, every retained section cocycle restricts to a
homomorphism on \(K\), since its actual geometric action there is
trivial. The image of the pro-two closure of \(k\) in an odd-order
geometric group is therefore the identity. Its image in the
two-primary group is geometric; every retained actual arithmetic
action is fixed. Additional odd-degree Galois fields over \(F\)
are also fixed, since \(k\) belongs to \(G_F\) and its closure
is pro-two. Thus powers of \(k\) give precisely permitted
conjugations and map onto every cyclic subquotient of \(R/F\).
This argument applies to arbitrary finite section quotients; it
does not assume a rational-point section or a class-two exponent
bound for their geometric group.

Put \(C=L\cap R\). Both fields are Galois over \(\mathbb Q\),
so \(C/F\) is cyclic of order \(2^s\) for some \(s\ge2\).
The lower bound follows from
\[
          F(B^{1/8})\subset C,\qquad
          [F(B^{1/8}):F]=4.
 \tag{E.170.6}
\]
The degree is a consequence of (E.170.1): the intermediate
radical subfields have the indicated degrees in the cyclic tower.
Write \(A=\operatorname{Gal}(C/F)\) additively.
The restriction of \(g\) to \(F\) fixes \(\sqrt B\), since it
fixes the marked eighth root. The polynomial
\(X^{2^{H-1}}-\sqrt B\) is irreducible over \(F\), by
(E.170.1). Therefore the restriction \(g|_F\) extends to an
automorphism \(t\) of \(R\) fixing \(B^{1/2^H}\): apply
\(g\) to coefficients and fix that chosen root. Both \(t\)
and \(g\) fix \(B^{1/8}\), so their discrepancy on \(C\) is
\[
             \delta=(g|_C)(t|_C)^{-1}\in4A.
 \tag{E.170.7}
\]
The group \(4A\) is exactly the kernel of restriction to the
degree-four field in (E.170.6).

Conjugation by an element with image \(a\in A\) changes this
discrepancy by \((1-\chi_2(g))a\). The arithmetic conjugation
action on a Kummer cyclic group is cyclotomic, independently of
the affine translation part. By (E.170.2),
\[
                 (1-\chi_2(g))A=4A.
 \tag{E.170.8}
\]
Powers of \(k\) surject onto \(A\), by (E.170.5). Choose one
whose image cancels \(\delta\). It fixes \(F\), all retained
arithmetic data and every odd geometric component. The conjugated
automorphism of \(L\) now agrees with \(t\) on \(C\), hence
the two extend to a single automorphism of \(LR\). Its
two-primary section value remains in exactly the same geometric
conjugacy class, and its marked eighth root remains fixed: the
extra Kummer translation is divisible by eight.

Chebotarev \citep[Theorem~8.31, p.~147]{MilneANT2020} gives
infinitely many unramified primes with this conjugacy class.
Being conjugate to a root-fixing automorphism means fixing some
root of the same polynomial, which reduces to the required root
in the prime field. We may exclude all primes dividing its
discriminant, the numerator or denominator of \(B\), and any
additional finite set. The full arithmetic and section conditions
are retained in the asserted conjugacy-invariant form.

For E.154, \(B=41V^8/(64W^8)\) has valuation one at 41. Its
quadratic squareclass is 41, and the pure-radical argument in that
result proves (E.170.1) for a cyclotomic field containing the
relevant roots of unity and \(\sqrt{41}\). The cusp-zero second
coordinate is literally zero in its marked Kummer model, so its
Frobenius fixes the eighth root. For the inverse use the inverse
compatible root. For the infinity direction use the displayed
Kummer combination: division by a chosen root of \(-1\),
rather than multiplication by it, gives exactly the rational
element \(-b/a\). The unit-one boundary relation is
\(\kappa(b/c)-\kappa(a/c)=\kappa(-1)\), so this marked
coordinate vanishes there. Its odd private valuation at 23 or
41 supplies (E.170.5), and the negative radical degree is the
one already proved in E.154. This verifies the stated applications.

The finite certificate \texttt{code/check\_two\_primary\_root.py}
checks the conjugation image in (E.170.8) and the exact vanishing
condition (E.170.7) for all affine parameters through modulus
\(2^9\). It also checks that retaining only a fourth-root mark
would be insufficient. The general field and conjugation
arguments are the proof above, not consequences of finite sampling.
\end{proof}
